
\begin{corollary}[Known after Gabber; a generic-field replacement without excellence]
\label{A:generic-plus-dimension}
Let $X$ be a finite-dimensional noetherian integral scheme, let $\ell$ be invertible on $X$, and put $K=\Frac(X)$. If $\ell=2$, assume that no residue field of $X$ is orderable. Then
\[
 \cd_\ell(X)\leq\cd_\ell(K)+\dim X\leq2\cd_\ell(K).
\]
The first estimate is sharp even for excellent regular affine schemes.
\end{corollary}
\begin{proof}
Theorem~\ref{A:strict-field} supplies, without excellence, the local fraction-field hypothesis in \citet[Exposé X, Corollaire 4.4]{A_SGA4}. That corollary gives the first inequality. Its proof applies the local specialization lower bound and the dimension-function criterion; no affine hypothesis is used. The second inequality is \citet[Exposé X, Corollaire 2.5]{A_SGA4}. The excellent Dedekind scheme of Theorem~\ref{A:nagata} has dimensions $3$, $2$, and $1$, respectively, and therefore realizes equality in the first bound.
\end{proof}

\subsection{A.2: a proved eventual-equality case beyond separably closed fields}

\begin{lemma}[Proved; quadratic descent with a finite dimension bound]
\label{A:quadratic-descent}
Let $f:T'\to T$ be a torsor under the constant group $\Z/2$, and suppose $\cd_2(T)<\infty$. Then $\cd_2(T)=\cd_2(T')$.
\end{lemma}
\begin{proof}
Finite étale pushforward is exact, so $\cd_2(T')\leq\cd_2(T)$. Put $a=\cd_2(T')$. For every sheaf $F$ of $\F_2$-vector spaces on $T$ there is an exact sequence
\[
 0\longrightarrow F\longrightarrow f_*f^*F
   \longrightarrow F\longrightarrow0.
\]
The first map is the diagonal and the second is the sum of the two coordinates, after an étale trivialization of the torsor. The sequence therefore descends. For $q>a$, its cohomology sequence gives
$H^q(T,F)\simeq H^{q+1}(T,F)$. Since the latter groups vanish in sufficiently high degrees by hypothesis, they vanish for every $q>a$. Reduction to $\F_2$ coefficients proves $\cd_2(T)\leq a$.
\end{proof}

\begin{theorem}[Proved under stated hypotheses; A.2 over fields]
\label{A:eventual-field}
Let $X$ be an integral affine scheme essentially of finite type over a field $k$, and let $\ell\neq\operatorname{char}k$. Suppose $n=\cd_\ell(\Frac X)<\infty$. Every nonempty open $W$ of $X$ contains a nonempty affine open $U$ such that
\[
 \cd_\ell(U')=n
 \quad\hbox{for every nonempty affine open }U'\subset U.
\]
If $\cd_\ell(k)<\infty$, the upper bound $\cd_\ell(V)\leq n$ holds for every affine open $V$ of $X$, before shrinking.
\end{theorem}
\begin{proof}
First suppose $X$ is of finite type, and put $d=\trdeg_k\Frac X$. If $\cd_\ell(k)<\infty$, Artin's affine theorem \citep[Exposé XIV, Corollaire 3.2]{A_SGA4}, followed by Hochschild--Serre descent, gives
\[
 \cd_\ell(V)\leq d+\cd_\ell(k)=n
\]
for every affine open $V$. The equality on the right is \citet[Exposé X, Théorème 2.1]{A_SGA4}.

That same field theorem says that the only remaining possibility is $\ell=2$, $k$ orderable, and $K=\Frac X$ nonorderable. In this case $\cd_2(k(i))<\infty$ and
\[
 n=d+\cd_2(k(i)).
\]
Write $-1$ as a finite sum of squares in $K$ and invert a common nonzero denominator. On the resulting nonempty open $X_1$, the identity remains valid in every residue field, so none is orderable. For every affine open $V\subset X_1$, Corollary~\ref{A:generic-plus-dimension} gives finite $\cd_2(V)$. The quadratic étale cover $V_{k(i)}\to V$ and Lemma~\ref{A:quadratic-descent} now give
\[
 \cd_2(V)=\cd_2(V_{k(i)})\leq d+\cd_2(k(i))=n.
\]
This also covers the case that the quadratic cover is split.

For the essentially finite-type version, write the coordinate ring as a localization of a finitely generated integral $k$-algebra. Any affine open and any constructible sheaf on it descend to an affine open at a finite localization stage. In the orderable case first invert the finitely many denominators in the displayed sum of squares. The upper bounds just proved hold at every such stage; continuity \citep[Exposé VII, Théorème 5.7]{A_SGA4} carries them to the localization. Finally apply Proposition~\ref{A:spread} inside $W$, intersecting with $X_1$ when necessary. Its one surviving generic class gives the opposite inequality simultaneously on every smaller nonempty affine open.
\end{proof}

\begin{proposition}[Known with derivation; the arithmetic test in the earlier study]
\label{A:arithmetic-test}
Let $K$ be a totally imaginary number field and $X=\Spec\mathcal O_K$. For every prime $\ell$,
\[
 \cd_\ell(X)=3,\qquad
 \cd_\ell(U)=2
 \quad (\,\varnothing\neq U\subsetneq X\text{ open}\,).
\]
Consequently the study is right that $X$ is not a counterexample satisfying the invertibility hypothesis of A.6. Its conclusion that A.6 remains undecided is superseded by Corollary~\ref{A:nagata-global-consequences}.
\end{proposition}
\begin{proof}
We use Artin--Verdier duality in the form \citet[Chapter II, Theorem 3.1, Corollary 3.2, and Proposition 2.3(a)]{A_Milne2006}, for arbitrary finite constructible étale sheaves. Since $K$ has no real places, no modified real-place term occurs. Duality gives vanishing above degree three. If $j:X[1/\ell]\to X$ and $F=j_!\mu_\ell$, then
\[
 H^3(X,F)^\vee
 =\Hom_X(F,\Gm)
 =\Hom_{X[1/\ell]}(\mu_\ell,\Gm)\neq0,
\]
so the bound three is attained.

For a proper nonempty open $U$, the localization sequence is
\[
 \bigoplus_{v\notin U}H^2(K_v,F)\longrightarrow H_c^3(U,F)
 \longrightarrow H^3(U,F)\longrightarrow0.
\]
Global and local duality identify the dual of the first arrow with restriction
\[
 \Hom_U(F,\Gm)\longrightarrow
 \prod_{v\notin U}\Hom_{K_v}(F,\Gm).
\]
It is injective: a morphism to $\Gm$ which vanishes at the generic point vanishes on the normal integral scheme $U$, and extension from $K$ to any $K_v$ is faithful on morphisms of finite generic sheaves. There is at least one omitted finite place. Hence $H^3(U,F)=0$, and continuity and coefficient dévissage give $\cd_\ell(U)\leq2$ for all torsion sheaves.

For completeness, the value two is attained without assuming that a generic class automatically extends to all of $U$. The finite group $H^1(U,\F_\ell)$ has some dimension $b$ by the same arithmetic finiteness theorem. Choose $b+1$ distinct closed points, put $Z$ equal to their union, and let $j:U-Z\to U$. From
\[
 0\longrightarrow j_!\F_\ell\longrightarrow\F_\ell
 \longrightarrow(\F_\ell)_Z\longrightarrow0
\]
we obtain a nonzero cokernel of
$H^1(U,\F_\ell)\to H^1(Z,\F_\ell)\simeq\F_\ell^{b+1}$
inside $H^2(U,j_!\F_\ell)$. Thus $\cd_\ell(U)=2$.
Every rational prime fails to be a unit on $X$ itself. On $X[1/\ell]$, the number $n$ in A.6 is indeed two: the generic term is two and each closed-point term is $1+1$.
\end{proof}

